\documentclass[10pt]{article}
\usepackage[margin=20mm]{geometry}
\usepackage{amsmath,amssymb,mathtools,fontspec,unicode-math,xurl,hyperref}
\IfFontExistsTF{DejaVu Serif}{\setmainfont{DejaVu Serif}}{\setmainfont{Latin Modern Roman}}
\setmathfont{latinmodern-math.otf}
\setlength{\parindent}{0pt}\setlength{\parskip}{6pt}\sloppy
\title{Initial scientific review of
Sparse equidistribution of geometric invariants of real quadratic fields}
\author{AI-assisted technical review}
\date{}
\begin{document}
\maketitle
Paper: M25-0315. Authors: Peter Humphries and Asbjørn Christian Nordentoft. DOI: https://doi.org/10.4171/JEMS/1691. Reviewed object: publisher Version of Record, JEMS 28 (2026), 5113–5192, 80 PDF pages; publisher publication date 25 August 2025. The frozen 2025 assignment is retained. This is not a review of an assumed equivalent arXiv revision.\par
\section*{Assessment and scope}
Assessment: Major, confined to a demonstrated failure of the stated regularisation and its continuous-spectrum proof dependency. The main equidistribution conclusions have not been disproved. One confirmed problem is established below by an admissible geodesic and the paper’s own Fourier expansion. The proposed local repair is natural, but its ease does not make the literal undefined integral a valid proof.\par
All 80 pages of extracted text were read; the extraction has defective mathematical glyphs, so 28 selected source pages were also visually inspected, including every page used for the confirmed problem and its immediate dependency. This is a substantive initial review, not a formal verification of every cited analytic theorem. Decision and adjudication remain pending; there are no human reviews or independent-acceptance claims.\par
\section*{Quantifiers and positive checks}
The distinctions among the four main theorems matter. Theorem 1.1 assumes a positive squarefree fundamental discriminant and a split odd prime level, with \(q\le D^\delta\), \(\delta<1/36\); its conditional range is \(\delta<1/12\). Theorem 1.2 fixes the level and allows a coset of a subgroup of relative size \(\gg D^{-\delta}\), with \(\delta<1/2826\) (conditionally \(1/4\)). Its fundamental discriminant need not be squarefree. Theorems 1.3–1.4 return to squarefree fundamental discriminants and level one; the discrepancy supremum is explicitly over injective balls.\par
The geometric volume lower bound in Proposition 3.9 is used together with the narrow class number formula. The trace estimate after the parabolic-block decomposition on p. 5138 yields \(\operatorname{vol}(\mathcal F_A(q))\gg\log\epsilon_D/\log(qD)\), rather than an unjustified bound from ordinary word length alone. The homological boundary contributions in Section 5 and Proposition 6.1 cannot generally be omitted when the modular curve has positive genus.\par
The stated level error \(q^{3+\varepsilon}D^{-1/12+\varepsilon}\) indeed gives the strict range \(\delta<1/36\). The fixed-level subgroup argument is also arithmetically consistent: the three squared Weyl-sum bounds give savings \(1/2826\), \(1/2114\), and \(1/1889\) after taking square roots. Character orthogonality supplies \((|H|/h_D^+)|H^\perp|=1\); it introduces no additional class-number factor. The slowest saving gives the printed \(1/2826\). These calculations do not extend the fixed-level theorem to growing levels.\par
For small scales, the three relative errors in (9.11) match at \(R=D^{-1/12}\) and \(D^{-5/12}\). At \(R=D^{-\delta}\) the final branch tends to zero for \(\delta<1/2\). Multiplication by the ball area, of order \(R^2\), is consistent with the discrepancy exponent. These exponent checks are conditional on the analytic identities and moment inputs, including the regularisation discussed next.\par
\section*{Confirmed problem}
Problem 1 (Major: the cusp cutoffs do not define the claimed regularised integral).  Location. Section 5.3, p. 5154 (physical PDF p. 42), (5.8)–(5.11), especially Lemma 5.4. The dependency is explicit in Lemma 6.10 and its proof, pp. 5165–5166 (physical pp. 53–54).\par
The definition is \[\mathfrak F_b(Y)=
 \{z\in\mathbb H:0<\Re(\sigma_b^{-1}z)<1,\
                         \Im(\sigma_b^{-1}z)\ge Y\},
 \qquad \sigma_b\in\mathrm{SL}_2(\mathbb Z),\quad\sigma_b\infty=b.\] It is asserted to be independent of \(\sigma_b\), and equivariant under the modular group. Both assertions fail for these subsets of \(\mathbb H\). For \(b=\infty\), the allowed matrices \(I\) and \(T=\left(\begin{smallmatrix}1&1\\0&1\end{smallmatrix}\right)\) give strips \(0<x<1\) and \(1<x<2\), respectively. The point \(1/2+2iY\) belongs only to the first. With the first choice fixed, \(T\mathfrak F_\infty(Y)\ne\mathfrak F_{T\infty}(Y)\).\par
There is a direct failure of Lemma 5.4, not merely a set-theoretic ambiguity. Take the permitted geodesic \[\mathcal C=\{iy:y>0\},\qquad b_1=\infty,\quad b_2=0,\qquad
 \sigma_\infty=I,\quad
 \sigma_0=\begin{pmatrix}0&-1\\1&0\end{pmatrix}.\] For \(z=iy\), both \(\sigma_\infty^{-1}z=iy\) and \(\sigma_0^{-1}z=i/y\) have real part zero. Consequently \[\mathcal C\setminus
       (\mathfrak F_\infty(Y)\cup\mathfrak F_0(Y))=\mathcal C
 \quad\hbox{for every }Y.\] The displayed limit (5.11) has not truncated either endpoint. Boundary sets of area zero cannot be discarded here: the whole integration curve lies on the excluded boundaries.\par
To verify divergence independently, put \(s=1/2+i\). The Fourier expansion (4.10), p. 5142, gives \[E(iy,s)=y^s+\varphi(s)y^{1-s}+O(e^{-c y}\,\mathrm{poly}(y)),
 \qquad
 \varphi(s)=\frac{\xi(1-2i)}{\xi(1+2i)}.\] On this axis, the \(x\) derivative vanishes by the symmetry of the Fourier coefficients, and \(R_0=iy\partial_x+y\partial_y\). Thus the upper-end integral from \(iA\) to \(iB\) equals \[i\bigl(B^s+\varphi(s)B^{1-s}\bigr)+C_A+o(1).\] Choose \(\theta\) with \(e^{i\theta}+\varphi(s)e^{-i\theta}\ne0\) and set \(B_n=e^{\theta+2\pi n}\). The displayed leading term is a fixed nonzero complex number times \(B_n^{1/2}\), so the ordinary improper integral does not converge. The endpoints, scaling matrices and spectral parameter satisfy the lemma’s stated hypotheses.\par
\subsection*{Impact, contrary evidence, and repair boundary}
Lemma 6.10 uses exactly these cusp zones in its Stokes truncation and invokes Lemma 5.4 for the Eisenstein boundary periods. Hence the literal definitions do not supply the claimed continuous-spectrum identities used subsequently in the level and subgroup arguments. This is a failure in a core proof dependency, which warrants Major at the proof level; it is not evidence that the equidistribution assertions themselves are false. The cuspidal decay argument and the elementary exponent checks above are not refuted by this example.\par
A precise candidate correction is to use full horoballs \[\mathcal H_b(Y)=
       \{z:\Im(\sigma_b^{-1}z)\ge Y\}.\] Two allowed scaling matrices differ by right multiplication by \(\pm T^n\), which preserves imaginary part. This restores scaling independence and modular equivariance. In the example it leaves the finite segment \(z=iy\), \(1/Y<y<Y\). Indeed modular invariance gives \(E(iY,s)=E(i/Y,s)\), so the corrected symmetrically truncated integral is exactly zero. This supplies genuine counterevidence against claiming that the intended regularised value is impossible.\par
A quotient-space reading can likewise be made meaningful only by taking the complete preimage of the cusp neighbourhood when truncating an arbitrary lifted geodesic. Changing the strict strip boundaries alone does not restore scaling independence. The author should replace the cutoff definition and check the general two-endpoint cancellation, the scaling law (5.10), and the Stokes boundary terms with that same convention. This review verifies the correction for the exhibited case, not every subsequent regularised identity.\par
\section*{Unresolved checks and limitations}
The topological terms in (6.19)–(6.20), p. 5168, appear without the inverse Laplace eigenvalue present in Lemma 6.10. The original page was checked, but this review has not supplied a permitted instance with a provably nonzero total topological term. It is therefore recorded as an unresolved normalisation check, not a second confirmed error.\par
Similarly, the equality between a periodised ball kernel and an indicator in Section 9 has an injectivity qualification. Extending the proof uniformly through larger balls or elliptic centres requires care with multiplicities. This review has not established a counterexample to Theorem 1.3 from that issue.\par
The adèlic normalisation chain, raising operators, parity factor, and fixed-level character counting were traced. The external toric-period theorem with its modified archimedean test vector, the general fractional-moment extension in (9.6), and all uniform analytic estimates were not independently reproved. No claim of exhaustive mathematical certification is made.\par
\section*{Supporting checks}
The acquired PDF has SHA256\par
00000TOKEN\par
and its manifest, title encoding, author names, issue year and publisher date are retained. A bounded exact-title search and the author’s research page located an official author-maintained Errata.pdf. Its four pages were read and viewed; the retrieved copy contains no entry for this paper. Its absence of an entry is not a claim that no correction exists elsewhere. No separate supplementary dataset or author code was located, and none was executed.\par
The current web initial-review executor performed the source reading, analysis, report preparation and actual source/report image inspections. An existing native auxiliary independently checked only the supplied exponent algebra, coset involution and cusp-counterexample calculation; it did not inspect the paper, run author code, or perform acceptance. Its role is recorded separately. The controller’s GPT-6/Pro UI observations are attributed to the controller; backend identity is not independently asserted.\par
Actual Tectonic 0.17.0 compilation, exit status, complete logs, dependency snapshots and retained payloads, exact input/output hashes, PDF parsing, and separate visual-inspection receipts accompany this report. Render-time false inspection records remain unchanged. This report is an initial-review deliverable with scientific decision pending.\par
\end{document}
